\documentclass[12pt]{article}
\usepackage[utf8]{inputenc}
\usepackage[T1]{fontenc}
\usepackage{hyperref}
\title{Mayi-Thakurti Language Resources for K-12 Curriculum Development}
\author{Compiled by Claude Sonnet 5.5 (claude-sonnet-5-5) for Andr\'as Kornai}
\date{}
\begin{document}
\maketitle

\section{Introduction}
Mayi-Thakurti (ISO 639-3: xyt; AIATSIS code G16) is an extinct Mayi language of the Gulf Country of north-west Queensland. Breen counts it in a dialect chain with Ngawun, Mayi-Kulan, Mayi-Yapi and Wunumara. Gerhardt Laves recorded material on it, which Breen describes as making Mayi-Thakurti the second best known of the Mayi group. The Mayi group as a whole is thinly documented.

\section{Primary Educational Resources}
No primers, textbooks or school programmes were found.

\section{Linguistic and Reference Materials}
\begin{itemize}
\item Gavan Breen, review of Laves's Mayi-Thakurti material (PDF, ANU). This is the most substantial item found specific to this language. \url{https://www0.anu.edu.au/linguistics/nash/aust/laves/Breen-MT.pdf}
\item Gavan Breen (1981), \emph{The Mayi Languages of the Queensland Gulf Country} (Glottolog reference record). \url{https://glottolog.org/resource/reference/id/64840}
\item AIATSIS Austlang entry G16 for the Mayi-Thakurti area. \url{https://collection.aiatsis.gov.au/austlang/language/G16}
\item Glottolog record for ISO code mnt, which groups the Mayi varieties. \url{https://glottolog.org/resource/languoid/iso/mnt}
\end{itemize}

\section{Religious and Community Resources}
None found.

\section{Digital and Audio Resources}
None found. The Laves material is archival, and access arrangements were not established in the searches.

\section{Curriculum Development Context}
Of the Mayi varieties, Mayi-Thakurti probably has the most material to draw on, through Laves's recordings and Breen's analysis. Any use should be led by descendant communities and arranged with the holding archives. Without a living speaker community or existing programmes, a heritage curriculum would be mostly about the language and its region, with a small amount of recorded vocabulary.
\end{document}